\documentclass[aspectratio=169]{beamer}
\usepackage{iftex}
\ifPDFTeX\usepackage[T1]{fontenc}\usepackage{lmodern}\fi
\usepackage[czech]{babel}\usepackage{listings}
\newcommand{\CVUTTemplatePath}{../../../utils/presentation-template}
\usepackage{\CVUTTemplatePath/beamerthemeCVUT}
\setbeamertemplate{footline}[frame number]\beamertemplatenavigationsymbolsempty
\newcommand{\LANG}{CZ}\newcommand{\FRONTPAGE}{dark}
\title{Analýza výstupů programu (profiling)}
\subtitle{Základy programování\\Matouš Cejnek\\U12110, FS, ČVUT v Praze}
\author{}\institute{}\date{}
\renewcommand{\headlineA}{}\renewcommand{\headlineB}{}\renewcommand{\headlineC}{}
\definecolor{CodeBlue}{RGB}{38,83,125}\definecolor{SoftGray}{RGB}{242,244,247}
\definecolor{CodeGreen}{RGB}{38,125,94}\definecolor{WarningRed}{RGB}{170,45,45}
\setbeamerfont{normal text}{size=\large}\setbeamerfont{frametitle}{size=\LARGE}
\setbeamercolor{block title}{fg=white,bg=CodeBlue}\setbeamercolor{block body}{fg=black,bg=SoftGray}
\newcommand{\term}[1]{\textcolor{CodeBlue}{\textbf{#1}}}
\lstdefinestyle{python}{language=Python,basicstyle=\ttfamily\normalsize,
 keywordstyle=\color{CodeBlue}\bfseries,stringstyle=\color{CodeGreen},
 commentstyle=\color{gray},showstringspaces=false,columns=fullflexible,keepspaces=true}
\begin{document}
\begin{frame}\clearpage\thispagestyle{empty}\titlepage\end{frame}

\begin{frame}{Co si dnes odnesete?}
\begin{itemize}
\item Rozlišit pozorování, benchmark a profilování.
\item Připravit opakovatelné měření výkonu.
\item Použít \texttt{timeit}, \texttt{cProfile} a \texttt{pstats}.
\item Najít úzké hrdlo podle naměřených dat.
\item Ověřit, že optimalizace zachovala správnost.
\end{itemize}
\end{frame}

\begin{frame}{Proč program působí pomalu?}
\begin{itemize}
\item Výpočet provádí příliš mnoho práce.
\item Čeká na disk, síť, uživatele nebo jiný proces.
\item Opakovaně převádí či kopíruje stejná data.
\item Většina času může být ukryta v jediné funkci.
\end{itemize}
\begin{block}{Zásada}Nejprve měřit, potom měnit.\end{block}
\end{frame}

\begin{frame}{Tři odlišné otázky}
\begin{description}
\item[Měření] Jak dlouho trval tento běh?
\item[Benchmark] Která varianta je za daných podmínek rychlejší?
\item[Profilování] Ve kterých částech programu se spotřeboval čas?
\end{description}
\end{frame}

\begin{frame}{Výkon není jen čas}
\begin{itemize}
\item Doba odezvy: jak dlouho čeká uživatel.
\item Propustnost: kolik úloh program zpracuje za čas.
\item Paměť: kolik prostoru program potřebuje.
\item I/O: kolik dat čte, zapisuje nebo přenáší.
\item Vždy měříme veličinu důležitou pro konkrétní použití.
\end{itemize}
\end{frame}

\begin{frame}{Co ovlivňuje naměřený čas?}
\begin{itemize}
\item Velikost a tvar vstupních dat.
\item Zahřívání cache a předchozí běhy.
\item Zatížení operačního systému a jiných procesů.
\item Verze Pythonu, knihoven a hardware.
\item Náhodnost, síť a stav souborového systému.
\end{itemize}
\end{frame}

\begin{frame}{Korektní experiment}
\begin{enumerate}
\item Formulujeme konkrétní otázku a metriku.
\item Připravíme reprezentativní vstup.
\item Měníme jednu věc a ostatní podmínky držíme stejné.
\item Měření opakujeme a uchováme více výsledků.
\item Výsledek popíšeme spolu s podmínkami.
\end{enumerate}
\end{frame}

\begin{frame}[fragile]{Hrubé měření celého úseku}
\begin{lstlisting}[style=python]
from time import perf_counter

start = perf_counter()
result = solve(data)
elapsed = perf_counter() - start
print(f"{elapsed:.3f} s")
\end{lstlisting}
\begin{block}{\texttt{perf\_counter}}Monotónní hodiny s vysokým rozlišením pro měření intervalu.\end{block}
\end{frame}

\begin{frame}{Jeden běh může klamat}
\begin{itemize}
\item Malý rozdíl může být pouze šum měření.
\item První běh může načítat data a vytvářet cache.
\item Extrémní hodnota může zkreslit průměr.
\item Sledujeme více opakování, často minimum nebo medián.
\end{itemize}
\end{frame}

\begin{frame}[fragile]{\texttt{timeit}: malé opakované měření}
\begin{lstlisting}[style=python]
from timeit import repeat

times = repeat(
    "sum(values)",
    setup="values = list(range(1000))",
    number=10_000,
    repeat=5,
)
print(min(times))
\end{lstlisting}
\end{frame}

\begin{frame}{Co z měření vyjmout?}
\begin{itemize}
\item Přípravu dat oddělíme, pokud není součástí otázky.
\item Výpis do terminálu bývá pomalejší než vlastní výpočet.
\item Naopak načtení souboru ponecháme, měříme-li celý pracovní postup.
\item Hranice měřeného úseku musí odpovídat rozhodnutí, které děláme.
\end{itemize}
\end{frame}

\begin{frame}{Asymptotická složitost a skutečný čas}
\begin{itemize}
\item Big O popisuje růst nákladů s velikostí vstupu.
\item Neříká přesný čas ani konstantní režii.
\item Pomáhá předvídat chování pro velká data.
\item Profilování ukazuje konkrétní běh na konkrétním systému.
\end{itemize}
\end{frame}

\begin{frame}{Kdy je optimalizace důležitá?}
\begin{itemize}
\item Když měření ukáže problém vzhledem k požadavku.
\item Když program blokuje uživatele nebo nezvládá objem dat.
\item Když cena prostředků nebo energie není přijatelná.
\item Neoptimalizujeme automaticky každý krátký a čitelný úsek.
\end{itemize}
\end{frame}

\begin{frame}{Profil ukazuje rozdělení práce}
\begin{itemize}
\item Kolikrát byla funkce zavolána.
\item Kolik času strávila ve vlastním těle.
\item Kolik času spotřebovala včetně volaných funkcí.
\item Které volání přispívá k celkové době nejvíce.
\end{itemize}
\end{frame}

\begin{frame}[fragile]{Spuštění programu přes \texttt{cProfile}}
\begin{lstlisting}[basicstyle=\ttfamily\large]
python -m cProfile -o profile.prof app.py
\end{lstlisting}
\begin{itemize}
\item Program se spustí běžným způsobem.
\item Profil se uloží do souboru pro další analýzu.
\item Profilování přidává režii; časy nejsou čistý benchmark.
\end{itemize}
\end{frame}

\begin{frame}{Jak číst tabulku profileru?}
\begin{description}
\item[\texttt{ncalls}] Počet volání funkce.
\item[\texttt{tottime}] Čas ve funkci bez podvolání.
\item[\texttt{cumtime}] Čas ve funkci včetně podvolání.
\item[\texttt{percall}] Přepočet času na jedno volání.
\end{description}
\end{frame}

\begin{frame}{Vlastní a kumulativní čas}
\begin{itemize}
\item Vysoký \texttt{tottime}: drahá práce přímo ve funkci.
\item Vysoký \texttt{cumtime}, nízký \texttt{tottime}: drahé podvolání.
\item Velký počet levných volání může být významný dohromady.
\item Řazení podle různých sloupců odpovídá různým otázkám.
\end{itemize}
\end{frame}

\begin{frame}[fragile]{Třídění statistik}
\begin{lstlisting}[style=python]
from pstats import Stats

(
    Stats("profile.prof")
    .strip_dirs()
    .sort_stats("cumulative")
    .print_stats(15)
)
\end{lstlisting}
\end{frame}

\begin{frame}{Úzké hrdlo}
\begin{itemize}
\item Část systému, která omezuje výkon celku.
\item Nemusí to být nejdelší funkce zdrojového kódu.
\item Po optimalizaci se úzké hrdlo může přesunout jinam.
\item Zaměřujeme se na významný podíl celkového času.
\end{itemize}
\end{frame}

\begin{frame}{Typické zdroje zbytečné práce}
\begin{itemize}
\item Opakovaný výpočet stejné hodnoty uvnitř cyklu.
\item Lineární hledání tam, kde se hodí množina nebo slovník.
\item Zbytečné kopírování velkých kolekcí.
\item Mnoho malých zápisů místo dávkového zpracování.
\item Načítání stejných dat znovu a znovu.
\end{itemize}
\end{frame}

\begin{frame}{Optimalizační cyklus}
\begin{enumerate}
\item Stanovíme základní měření a testy správnosti.
\item Profilujeme reprezentativní scénář.
\item Zvolíme jedno významné úzké hrdlo.
\item Provedeme malou srozumitelnou změnu.
\item Znovu spustíme testy, benchmark i profil.
\end{enumerate}
\end{frame}

\begin{frame}{Rychlejší, ale stále správně}
\begin{itemize}
\item Optimalizace nesmí změnit požadovaný výsledek.
\item Automatické testy chrání běžné i hraniční případy.
\item U přibližných algoritmů měříme také kvalitu výsledku.
\item Porovnáváme výkon až po ověření korektnosti variant.
\end{itemize}
\end{frame}

\begin{frame}{Čitelnost je také náklad}
\begin{itemize}
\item Složitější kód se hůře udržuje a testuje.
\item Mikrooptimalizace může přinést zanedbatelnou úsporu.
\item Dobrou optimalizaci doložíme měřením a vysvětlením.
\item Často nejvíce pomůže vhodnější algoritmus nebo datová struktura.
\end{itemize}
\end{frame}

\begin{frame}{Časté chyby při profilování}
\begin{itemize}
\item Měření nereprezentativních nebo příliš malých dat.
\item Porovnání jediného běhu bez opakování.
\item Záměna profileru za přesný benchmark.
\item Optimalizace malé položky bez vlivu na celek.
\item Chybějící kontrola správnosti po změně.
\end{itemize}
\end{frame}

\begin{frame}{Čas procesoru a skutečný čas}
\begin{description}
\item[Wall time] Doba, kterou pozoruje uživatel od začátku do konce.
\item[CPU time] Doba, kdy procesor skutečně provádí instrukce programu.
\end{description}
\begin{itemize}
\item Při čekání na disk či síť roste wall time, ale CPU time téměř ne.
\item Volba metriky závisí na otázce, kterou řešíme.
\end{itemize}
\end{frame}

\begin{frame}{Výpočet versus čekání}
\begin{description}
\item[CPU-bound] Výkon omezuje množství výpočtů.
\item[I/O-bound] Výkon omezuje čekání na vstup či výstup.
\end{description}
\begin{itemize}
\item Rychlejší algoritmus pomůže hlavně výpočetně náročné úloze.
\item Dávkování a omezení čekání pomáhá u I/O operací.
\end{itemize}
\end{frame}

\begin{frame}{Profilujeme správný scénář?}
\begin{itemize}
\item Start aplikace a dlouhý provoz mohou mít jiná úzká hrdla.
\item Malý testovací soubor nemusí reprezentovat skutečná data.
\item Vzácná pomalá operace může být pro uživatele zásadní.
\item Profilujeme konkrétní pracovní postup, který chceme zlepšit.
\end{itemize}
\end{frame}

\begin{frame}{Jak výkon roste se vstupem?}
\begin{itemize}
\item Změříme několik velikostí vstupu, ne pouze jednu.
\item Sledujeme tvar růstu času a paměti.
\item Zdvojnásobení dat může odhalit kvadratický algoritmus.
\item Jednorázová režie dominuje hlavně u malých vstupů.
\end{itemize}
\end{frame}

\begin{frame}{Volací strom}
\begin{itemize}
\item Tabulka ukáže drahé funkce, strom ukáže cestu k nim.
\item Stejná funkce může být volána z několika míst.
\item Kontext volání pomáhá zvolit místo vhodné změny.
\item Optimalizace volajícího může odstranit mnoho podvolání najednou.
\end{itemize}
\end{frame}

\begin{frame}{Paměť a výkon}
\begin{itemize}
\item Velké kolekce zabírají paměť a zvyšují režii kopírování.
\item Generátor zpracovává hodnoty postupně místo uložení všech.
\item Cache šetří výpočet za cenu dodatečné paměti.
\item Sledujeme špičkovou spotřebu i životnost objektů.
\end{itemize}
\end{frame}

\begin{frame}{Výsledek měření}
\begin{itemize}
\item Uvedeme verzi kódu, prostředí a velikost vstupu.
\item Zapíšeme metriku, počet opakování a způsob agregace.
\item Porovnáme výchozí a upravenou variantu stejným postupem.
\item Zaznamenáme také dopad na čitelnost a spotřebu paměti.
\end{itemize}
\end{frame}

\begin{frame}{Shrnutí}
\begin{itemize}
\item Výkon vyhodnocujeme podle konkrétního požadavku a metriky.
\item \texttt{timeit} porovnává malé varianty opakovaným měřením.
\item \texttt{cProfile} a \texttt{pstats} ukazují rozdělení času mezi funkce.
\item Optimalizujeme naměřené úzké hrdlo, ne dojem ze zdrojového kódu.
\item Každou změnu ověříme testy a novým měřením.
\end{itemize}
\end{frame}
\end{document}
